\chapter{Conclusiones}
\label{chap:conclusiones}

\noindent Este capítulo presenta las conclusiones derivadas de la investigación, respondiendo al proyecto de tesis en su conjunto mediante la síntesis de los hallazgos principales, la confirmación de la hipótesis, las contribuciones metodológicas y científicas, así como las limitaciones identificadas y las líneas futuras de investigación.

% ============================================================================
% QUÉ SE PRESENTÓ (PROBLEMA)
% ============================================================================

Este estudio abordó el problema fundamental de la ausencia de un sistema que integre de forma sinérgica múltiples biomarcadores de vida libre en una clasificación del comportamiento sedentario que sea robusta, empíricamente validada e interpretable. Los análisis estadísticos univariados han resultado insuficientes para capturar las complejas interacciones no lineales entre la actividad física, la función autonómica y la respuesta metabólica, donde la relación entre biomarcadores como la variabilidad de la frecuencia cardíaca y la actividad física es intrínsecamente no lineal y dependiente del contexto, requiriendo enfoques de modelado que puedan manejar la incertidumbre y la vaguedad inherentes a los sistemas biológicos. Este vacío metodológico motivó el desarrollo de un clasificador objetivo, interpretable y validado para el comportamiento sedentario en condiciones de vida libre utilizando datos longitudinales recopilados a través de dispositivos de consumo, respondiendo a la pregunta de investigación: ¿De qué manera un sistema basado en lógica difusa puede clasificar el comportamiento sedentario a partir de datos biométricos recopilados en condiciones de vida libre por dispositivos ``\textit{wearables}''?

% ============================================================================
% MÉTODO PARA RESOLVER EL PROBLEMA
% ============================================================================

Para resolver este problema, esta investigación desarrolló un sistema de evaluación basado en un sistema de inferencia difusa tipo \textit{Mamdani} que combina descubrimiento de patrones data-driven mediante \textit{clustering} no supervisado con codificación de conocimiento experto mediante lógica difusa. El enfoque se apartó de una hipótesis inicial centrada en correlacionar datos objetivos con percepciones subjetivas, para en su lugar, validar el sistema experto contra una verdad operativa derivada de los propios datos mediante análisis de conglomerados no supervisado. La metodología comprendió la derivación de cuatro variables semanales compuestas (Actividad relativa, Superávit calórico basal, Delta cardíaco, \textit{HRV-SDNN}) mediante normalización antropométrica y agregación semanal sobre 9,185 días de registro longitudinal retrospectivo, resultando en 1,337 semanas válidas. El algoritmo \textit{K-Means} identificó dos perfiles de comportamiento con separación moderada (Silhouette = 0.232) que sirvieron como verdad operativa, mientras que el sistema difuso \textit{Mamdani} fue diseñado con 5 reglas lingüísticas fundamentadas fisiológicamente y funciones de pertenencia triangulares calibradas a partir de percentiles observados en los datos. La evaluación del desempeño se realizó mediante validación cruzada Leave-One-User-Out (\textit{LOUO}), y se examinó la contribución de los componentes del modelo a través de análisis de sensibilidad mediante ablación sistemática.

% ============================================================================
% SÍNTESIS DE RESULTADOS PRINCIPALES
% ============================================================================

Los resultados principales demuestran que el sistema de inferencia difusa alcanzó una alta concordancia (\textit{F1-Score} \textit{LOUO} = 0.780 $\pm$ 0.167) con la clasificación objetiva derivada del análisis de conglomerados, con precisión de 0.800 y sensibilidad (\textit{recall}) de 0.783. El desempeño se mantuvo consistente a través de 7 de 10 usuarios con \textit{F1-Score} $\geq$ 0.65, con variabilidad inter-sujeto de CV=21.4\%, indicativo de variabilidad moderada esperada en estudios de vida libre. El análisis de \textit{clustering} \textit{K-Means} ($k=2$) identificó dos perfiles de comportamiento con separación moderada (índice de Silhouette = 0.232) y distribución asimétrica (Clúster 0 [ACTIVO]: 402 semanas [30.1\%], Clúster 1 [SEDENTARIO]: 935 semanas [69.9\%]), donde los centroides exhibieron diferencias significativas en las cuatro variables derivadas ($p < 0.001$ en todas las comparaciones Mann-Whitney, con tamaños de efecto grandes: Cohen's d $>$ 0.80). El análisis de ablación reveló un fenómeno contraintuitivo de alta relevancia metodológica: la variabilidad de la frecuencia cardíaca (\textit{HRV-SDNN}) no demostró poder discriminatorio significativo en análisis univariado (prueba de Mann-Whitney: U = 184,180, $p=0.562$, Cohen's d = 0.051, efecto despreciable), sin embargo, al excluir \textit{HRV-SDNN} del sistema de inferencia difusa mediante ablación, el desempeño clasificatorio colapsó de \textit{F1-Score}=0.840 a \textit{F1-Score}=0.420, representando una pérdida del 50\% en capacidad predictiva, evidenciando que \textit{HRV-SDNN} actúa como modificador de efecto en interacción con otras variables biométricas.

% ============================================================================
% CONFIRMACIÓN DE HIPÓTESIS
% ============================================================================

La hipótesis conceptual establecía que la clasificación del comportamiento sedentario generada por el sistema de inferencia difusa, basado en reglas que integran biomarcadores de actividad y función cardiovascular, exhibe una alta concordancia con la clasificación objetiva obtenida mediante un análisis de conglomerados no supervisado sobre los mismos datos. Esta hipótesis fue confirmada con evidencia estadística sólida: el coeficiente Kappa de Cohen entre las clasificaciones del \textit{clustering} y del sistema difuso alcanzó 0.56 (concordancia moderada-sustancial), significativamente superior al acuerdo esperado por azar ($p < 0.001$). El \textit{F1-Score} de 0.780 en validación \textit{LOUO} indica que el modelo difuso logró replicar la estructura clasificatoria identificada por el \textit{clustering}, generalizando efectivamente a individuos no vistos durante el entrenamiento. Este resultado confirma la viabilidad de la arquitectura híbrida propuesta, donde el conocimiento empírico descubierto por métodos no supervisados puede ser codificado en reglas lingüísticas interpretables que preservan la capacidad predictiva, validando el enfoque de utilizar \textit{clustering} como verdad operativa en contextos donde estándares de oro clásicos no están disponibles.

% ============================================================================
% CONTRIBUCIONES DEL TRABAJO
% ============================================================================

Las contribuciones de este trabajo se manifiestan en múltiples niveles complementarios. Metodológicamente, se aporta una estrategia de validación robusta para sistemas de clasificación en salud cuando no se dispone de un estándar de oro clínico, basada en la convergencia de un enfoque empírico (descubrimiento de patrones no supervisado) y uno experto (reglas basadas en conocimiento). Esta arquitectura híbrida representa una contribución metodológica al campo de la salud digital, demostrando que es posible combinar ambos enfoques de manera sinérgica para lograr clasificación robusta y clínicamente interpretable. La ingeniería de las variables Actividad Relativa y Superávit Calórico Basal se presenta como una solución eficaz para la normalización fisiológica de datos de dispositivos portátiles, permitiendo comparaciones más equitativas entre individuos. Científicamente, se identificó y caracterizó un fenómeno de modificación de efecto donde la variabilidad de la frecuencia cardíaca actúa como modulador contextual crítico a pesar de carecer de asociación univariada significativa, aportando evidencia sobre la existencia de modificación de efecto y supresión estadística en fisiología del ejercicio que no había sido claramente demostrada en estudios de vida libre con \textit{wearables} de consumo. Técnicamente, se demostró que el paradigma \textit{Bring Your Own Device} no solo es viable metodológicamente, sino que puede generar datos de alta calidad analítica sin comprometer la integridad de las mediciones, representando un avance significativo en la democratización de la investigación digital en salud. La consolidación de una arquitectura computacional flexible que puede integrarse en sistemas de monitoreo remoto constituye una contribución tecnológica con potencial de transferencia hacia el diseño de intervenciones personalizadas, donde el código fuente documentado y los protocolos de preprocesamiento están disponibles para replicación en cohortes independientes.

% ============================================================================
% LIMITACIONES DEL ESTUDIO
% ============================================================================
\section{Limitaciones del Estudio}
\label{sec:limitaciones}

A pesar de los logros alcanzados, es importante reconocer limitaciones que contextualizan la interpretación de los resultados. El diseño observacional longitudinal retrospectivo no permite establecer relaciones causales, únicamente asociaciones temporales y predictivas, requiriéndose diseños experimentales controlados con asignación aleatoria para determinar causalidad. Aunque el \textit{clustering} \textit{K-Means} proporciona una clasificación objetiva y data-driven, no constituye un estándar de oro en el sentido clásico, limitando la validación de criterio del modelo debido a la ausencia de mediciones de referencia como calorimetría indirecta o actigrafía de investigación, limitación inherente a estudios de vida libre donde la obtención de estándares de oro es logísticamente inviable en seguimientos longitudinales de varios años. El paradigma \textit{BYOD} introduce heterogeneidad en las versiones de \textit{hardware} (\textit{Apple Watch} Series 3-9) y \textit{software} (watchOS 7-10) utilizadas por los participantes, donde diferencias menores en la precisión de los sensores podrían contribuir marginalmente a la variabilidad observada, aunque los algoritmos propietarios de Apple mantienen consistencia en las métricas reportadas.

La muestra de $N=10$ usuarios, aunque generó 1,337 semanas válidas suficientes para \textit{clustering} y validación estadística, limita la capacidad para realizar análisis de subgrupos y detectar efectos menores que requerirían mayor poder estadístico. El muestreo por conveniencia (usuarios de \textit{Apple Watch} con disposición a compartir datos) puede limitar la generalización de los resultados a poblaciones que no tienen acceso a dispositivos portátiles o que no mantienen uso consistente, mientras que la composición principalmente de adultos jóvenes universitarios (edad media: 25-35 años) limita la aplicabilidad a otros grupos etarios o poblaciones con comorbilidades crónicas. Los resultados provienen de una cohorte de Chihuahua, México, con características sociodemográficas, climáticas y culturales específicas, requiriendo cautela en la generalización a otras regiones con patrones de actividad física diferentes. Los datos longitudinales abarcan un período que incluye la pandemia COVID-19 (2020-2023), durante la cual se implementaron confinamientos y restricciones de movilidad que pudieron alterar artificialmente los patrones de actividad física, aunque la agregación semanal mitiga efectos de corto plazo. La dependencia exclusiva de \textit{Apple Watch} limita la generalización a usuarios de otras plataformas, mientras que el diseño retrospectivo no permitió implementar mediciones concurrentes de validación, requiriéndose estudios futuros con diseño prospectivo para subsanar esta limitación.

% ============================================================================
% TRABAJO FUTURO
% ============================================================================
\section{Trabajo Futuro}
\label{sec:trabajo_futuro}

Las líneas futuras de investigación incluyen validación prospectiva con estándares de oro mediante mediciones concurrentes de actigrafía de investigación y calorimetría indirecta para cuantificar la validez de criterio del modelo difuso versus técnicas de referencia. El escalamiento a cohortes heterogéneas con mayor diversidad etaria, socioeconómica y clínica, incluyendo poblaciones con enfermedades crónicas, permitiría evaluar la robustez del modelo en condiciones de alta heterogeneidad y extender su aplicabilidad a grupos poblacionales diversos. La implementación en intervenciones conductuales randomizadas mediante aplicaciones móviles que utilicen el modelo difuso para generar retroalimentación en tiempo real permitiría evaluar su efectividad en reducir el sedentarismo mediante ensayos clínicos controlados aleatorizados, validando no solo la capacidad clasificatoria del modelo sino también su impacto en cambios conductuales reales. La exploración de variabilidad temporal mediante análisis de patrones circadianos y semanales del comportamiento sedentario mediante descomposición de series temporales permitiría identificar momentos de mayor susceptibilidad a intervenciones conductuales, optimizando el timing de las intervenciones para maximizar su efectividad. La integración con biomarcadores inflamatorios y metabólicos mediante vinculación de los perfiles de sedentarismo identificados con mediciones sanguíneas de proteína C reactiva, glucosa, lípidos y citoquinas inflamatorias permitiría caracterizar las vías biológicas mediando la asociación sedentarismo-enfermedad, proporcionando evidencia mecanística sobre los procesos fisiopatológicos subyacentes. Finalmente, la adaptación del protocolo para dispositivos portátiles alternativos (Garmin, Fitbit, Samsung) mediante ajuste de las funciones de pertenencia a los centroides observados en nuevas cohortes permitiría extender la aplicabilidad del modelo más allá del ecosistema Apple, democratizando aún más el acceso a herramientas de evaluación del comportamiento sedentario basadas en datos de vida libre.
